% Concise source-only Overleaf report matched to the presentation-ready simulation.
% First run Three_Structures_LSTM_VSCode_Windows_Fixed_AutoPush_PresentationReady.py.
% Then copy PNG files from 01_Report_Figures/<timestamp>/ into figures/simulation/.

\documentclass[11pt]{article}
\usepackage[margin=0.9in]{geometry}
\usepackage{
amsmath,amssymb,graphicx,booktabs,array,tabularx,
hyperref,xcolor,float,enumitem,caption,subcaption
}
\usepackage[most]{tcolorbox}

\hypersetup{
colorlinks=true,
linkcolor=blue,
citecolor=blue,
urlcolor=blue
}

\graphicspath{{figures/}{figures/simulation/}}

\definecolor{graybg}{RGB}{247,247,247}
\definecolor{grayframe}{RGB}{90,90,90}
\definecolor{paperbg}{RGB}{244,238,255}
\definecolor{paperframe}{RGB}{103,71,170}
\definecolor{databg}{RGB}{232,248,239}
\definecolor{dataframe}{RGB}{31,122,65}
\definecolor{archbg}{RGB}{231,243,255}
\definecolor{archframe}{RGB}{23,88,156}
\definecolor{trainbg}{RGB}{255,248,220}
\definecolor{trainframe}{RGB}{176,126,0}
\definecolor{resultbg}{RGB}{255,239,239}
\definecolor{resultframe}{RGB}{160,45,45}

\newtcolorbox{summarybox}[1]{
enhanced,breakable,
colback=graybg,colframe=grayframe,
title={#1},fonttitle=\bfseries,
arc=2mm,boxrule=0.8pt
}

\newtcolorbox{paperbox}[1]{
enhanced,breakable,
colback=paperbg,colframe=paperframe,
title={#1},fonttitle=\bfseries,
arc=2mm,boxrule=0.8pt
}

\newtcolorbox{databox}[1]{
enhanced,breakable,
colback=databg,colframe=dataframe,
title={#1},fonttitle=\bfseries,
arc=2mm,boxrule=0.8pt
}

\newtcolorbox{archbox}[1]{
enhanced,breakable,
colback=archbg,colframe=archframe,
title={#1},fonttitle=\bfseries,
arc=2mm,boxrule=0.8pt
}

\newtcolorbox{trainbox}[1]{
enhanced,breakable,
colback=trainbg,colframe=trainframe,
title={#1},fonttitle=\bfseries,
arc=2mm,boxrule=0.8pt
}

\newtcolorbox{resultbox}[1]{
enhanced,breakable,
colback=resultbg,colframe=resultframe,
title={#1},fonttitle=\bfseries,
arc=2mm,boxrule=0.8pt
}

\newcommand{\simfigure}[4]{%
\begin{figure}[H]
\centering
\IfFileExists{figures/simulation/#2}{%
\includegraphics[width=#1\textwidth]{figures/simulation/#2}
}{%
\begin{tcolorbox}[
colback=white,
colframe=black!45,
title={#3},
width=0.92\textwidth
]
\centering
\vspace{0.18cm}
\textbf{Insert the figure generated by the Python code.}\\[0.12cm]
Expected filename:
\texttt{\detokenize{#2}}\\[0.12cm]
\small
Copy it from
\texttt{01\_Report\_Figures/<timestamp>/}
into
\texttt{figures/simulation/}.
\vspace{0.18cm}
\end{tcolorbox}
}
\caption{#4}
\end{figure}
}

\newcommand{\simpair}[8]{%
\begin{figure}[H]
\centering
\begin{subfigure}[t]{0.485\textwidth}
\centering
\IfFileExists{figures/simulation/#2}{%
\includegraphics[width=\linewidth]{figures/simulation/#2}
}{%
\begin{tcolorbox}[
colback=white,colframe=black!45,
title={#3},width=\linewidth
]
\centering
\scriptsize
Expected filename:\\
\texttt{\detokenize{#2}}
\vspace{0.15cm}
\end{tcolorbox}
}
\caption{#4}
\end{subfigure}
\hfill
\begin{subfigure}[t]{0.485\textwidth}
\centering
\IfFileExists{figures/simulation/#6}{%
\includegraphics[width=\linewidth]{figures/simulation/#6}
}{%
\begin{tcolorbox}[
colback=white,colframe=black!45,
title={#7},width=\linewidth
]
\centering
\scriptsize
Expected filename:\\
\texttt{\detokenize{#6}}
\vspace{0.15cm}
\end{tcolorbox}
}
\caption{#8}
\end{subfigure}
\caption{#5}
\end{figure}
}

\title{
Concise Report on Three-Structure LSTM-Based\\
Dynamic System Identification of Actuator Response
}
\author{Hussein Zolfaghari}
\date{\today}

\begin{document}
\maketitle

\begin{abstract}
This report summarizes a nonlinear dynamic-system identification study in
which coil current is used to predict actuator displacement and Lorentz force.
Three structures are implemented and compared: series, series--parallel, and
parallel/free-running prediction. Four DC-offset experiments are used for
training and validation, while the 147~mA experiment is kept completely unseen
for final testing.
\end{abstract}

\section{Purpose and Main Idea}

The known input vector is
\[
u(k)=
\begin{bmatrix}
I(k) & \Delta I(k) & I_{\mathrm{DC}}
\end{bmatrix}^{T},
\]
and the output vector is
\[
y(k)=
\begin{bmatrix}
x(k) & F(k)
\end{bmatrix}^{T},
\]
where $I$ is coil current, $x$ is displacement, and $F$ is Lorentz force.

A history of 120 samples is used to predict the next output. The three
implemented structures are

\begin{table}[H]
\centering
\caption{Implemented prediction structures}
\begin{tabularx}{\textwidth}{lX}
\toprule
Structure & Mathematical input--output relation \\
\midrule
Series &
$[I,\Delta I,I_{\mathrm{DC}}]_{\mathrm{past}}
\rightarrow[\hat{x}(k+1),\hat{F}(k+1)]$ \\[1mm]

Series--parallel &
$[I,\Delta I,I_{\mathrm{DC}},x,F]_{\mathrm{past}}
\rightarrow[\hat{x}(k+1),\hat{F}(k+1)]$ \\[1mm]

Parallel &
$[I,\Delta I,I_{\mathrm{DC}},\hat{x},\hat{F}]_{\mathrm{past}}
\rightarrow[\hat{x}(k+1),\hat{F}(k+1)]$ \\
\bottomrule
\end{tabularx}
\end{table}

\section{Connection to the Reference Papers}

\begin{paperbox}{Connection to Wang (2017)}
The implementation uses Wang's distinction between series--parallel
identification, where measured output history is supplied, and parallel
identification, where predicted output history is fed back. The present code
does not reproduce Wang's complete convex multi-LSTM ensemble; it is therefore
a Wang-inspired implementation.
\end{paperbox}

\begin{paperbox}{Connection to Ogunmolu et al. (2016)}
The model follows the general idea of using deep dynamic neural networks for
nonlinear input--output identification. A static nonlinear transformation of
the known inputs is followed by recurrent LSTM dynamics, giving a
Hammerstein-style neural architecture. The code is a PyTorch implementation,
not a direct conversion of the original FARNN source code.
\end{paperbox}

\section{Dataset and Data Usage}

The development datasets are

\[
67,\ 87,\ 107,\ 127~\mathrm{mA},
\]

and the independent test dataset is

\[
147~\mathrm{mA}.
\]

\begin{databox}{Cyclic blocked split across the chirp}
Because chirp frequency changes continuously, a single contiguous split can
make one subset contain mostly low-frequency behavior and another subset
contain mostly high-frequency behavior. The updated code therefore uses a
cyclic blocked split inside each development record:
\[
\text{train}\rightarrow\text{guard gap}\rightarrow
\text{validation}\rightarrow\text{guard gap}\rightarrow
\text{development test}\rightarrow\text{guard gap},
\]
and then repeats this pattern across the full chirp record. This makes the
training, validation, and internal development-test blocks cover different
frequency regions while reducing leakage from overlapping LSTM windows.
The complete 147~mA record is still reserved only for the final unseen test.
\end{databox}

The current implementation uses 1000 training samples, 500 validation samples,
500 internal development-test samples, and a 120-sample guard gap in each
cycle. Normalization statistics are calculated from training blocks only.

\simfigure{0.92}{00_cyclic_split_schematic.png}
{Cyclic split schematic}
{Schematic of the repeated train--validation--development-test block pattern
used across each chirp record.}

\simfigure{0.95}{02_data_usage_summary.png}
{Training, validation, development-test, and final-test data usage}
{Cyclic blocked split across the four development experiments. The complete
147~mA record is reserved for final testing.}

The code also produces colored current plots for every operating condition:

\begin{itemize}[nosep]
\item blue: training blocks;
\item orange: validation blocks;
\item purple: internal development-test blocks;
\item gray: guard-gap samples not used for optimization;
\item green: independent 147~mA final test record.
\end{itemize}

\section{Model Architecture and Training}

\begin{archbox}{Architecture summary}
The known inputs first pass through a static nonlinear block
\[
3\rightarrow24\rightarrow24
\]
with Tanh activation. The dynamic block is an LSTM with 64 hidden units.
The final prediction head is
\[
64\rightarrow32\rightarrow2.
\]
The two outputs are displacement and Lorentz force.
\end{archbox}

\begin{trainbox}{Training summary}
The series and series--parallel models are trained for 15 epochs using Adam,
a learning rate of $10^{-3}$, weight decay, gradient clipping, and validation
checkpoint selection. The series--parallel model receives measured output
history, with small training noise added for robustness. The parallel model
starts from the series--parallel weights and is fine-tuned for 8 epochs using
50-step recursive rollouts. The measured-feedback ratio decreases from 0.9 to
0.1.
\end{trainbox}

\section{Simulation Results}

The figures below are generated directly by the presentation-ready Python
simulation. All prediction results use the same completely unseen 147~mA chirp.
The first 120 samples are used only to initialize feedback-based evaluation and
are excluded from the reported accuracy metrics.

\subsection{Training and Validation Behavior}

\simfigure{0.88}{01_all_structures_cost.png}
{Training and validation cost}
{Normalized training and validation MSE for series, series--parallel, and
parallel models. The parallel objective is more difficult because it includes
recursive multi-step prediction.}

\subsection{Overall Comparison on the Unseen Test}

The new comparison figures contain both the full test record and a zoomed
high-frequency region. The zoom prevents important local differences from
being hidden by the full-record scale.

\simfigure{0.96}{03_displacement_all_structures.png}
{Displacement comparison for all structures}
{Measured displacement together with series, series--parallel, and parallel
predictions over the full unseen record and a high-frequency zoom.}

\simfigure{0.96}{03_force_all_structures.png}
{Lorentz-force comparison for all structures}
{Measured Lorentz force together with series, series--parallel, and parallel
predictions over the full unseen record and a high-frequency zoom.}

\subsection{Structure-by-Structure Tracking}

The following figures are the clearest way to compare the three strategies,
because each structure is plotted against the measured response without being
compressed by a strongly diverging curve from another model.

\simfigure{0.96}{06_series_measured_vs_predicted.png}
{Series measured versus predicted response}
{Series/input-only prediction on the unseen 147~mA chirp. This model is the
baseline because it uses known input history but no output history.}

\simfigure{0.96}{06_series_parallel_measured_vs_predicted.png}
{Series--parallel measured versus predicted response}
{Series--parallel one-step prediction using measured displacement and force
history. This figure demonstrates the expected accuracy advantage of measured
output feedback.}

\simfigure{0.96}{06_parallel_measured_vs_predicted.png}
{Parallel measured versus predicted response}
{Parallel/free-running prediction after initialization. Predicted displacement
and force are fed back recursively, so phase, amplitude, and bias errors can
accumulate.}

\begin{resultbox}{Main comparison}
The series model provides the input-only baseline. The series--parallel model
is expected to give the most accurate one-step prediction because measured past
outputs are supplied at every step. The parallel model is the strongest test of
autonomous simulation because it must continue using its own predicted output
history. Therefore, the principal advantage of the series--parallel strategy is
more stable learning and more accurate short-horizon tracking, while the
parallel result reveals long-horizon recursive stability.
\end{resultbox}

\subsection{Prediction Errors}

The error is defined as
\[
e(k)=y(k)-\hat{y}(k).
\]
Direct error plots are more meaningful than a crowded response plot when one
structure has much larger drift than the others.

\simpair
{0.48}
{04_displacement_errors_all_structures.png}
{Displacement error}
{Displacement error histories}
{Prediction errors for all three structures on the same unseen test record.}
{04_force_errors_all_structures.png}
{Force error}
{Lorentz-force error histories}

\subsection{Regression and Accuracy Metrics}

The regression figures plot predicted values against measured values. Ideal
prediction follows
\[
\hat{y}=y.
\]
A slope close to one, an intercept close to zero, and regression $R^2$ close to
one indicate strong agreement.

\simfigure{0.96}{08_all_structures_regression_comparison.png}
{Regression comparison}
{Measured-versus-predicted regression plots for displacement and Lorentz force
for series, series--parallel, and parallel models.}

The simulation also generates individual regression figures:

\begin{itemize}[nosep]
\item \texttt{07\_series\_regression.png};
\item \texttt{07\_series\_parallel\_regression.png};
\item \texttt{07\_parallel\_regression.png}.
\end{itemize}

\simfigure{0.90}{05_all_structures_accuracy_table.png}
{Accuracy table}
{RMSE, MAE, $R^2$, and system-identification fit percentage for displacement
and Lorentz force for all three structures.}

\begin{resultbox}{How to interpret the final result}
The advantage of series--parallel prediction should be demonstrated using the
measured-versus-predicted tracking figure, the error histories, the regression
plot, and the numerical table together. A visually good curve alone is not
sufficient. Similarly, weaker parallel performance should not be described as a
failure of one-step identification; it indicates that recursive free-running
simulation is a substantially harder task.
\end{resultbox}

\section{Generated Files}

Run the final simulation file first:
\begin{center}
\texttt{Three\_Structures\_LSTM\_VSCode\_Windows\_Fixed\_AutoPush\_PresentationReady.py}
\end{center}

After each successful run, the code creates

\begin{itemize}[nosep]
\item \texttt{01\_Report\_Figures/<timestamp>/};
\item \texttt{02\_Presentation\_Figures/<timestamp>/};
\item \texttt{03\_Complete\_Results/<timestamp>/}.
\end{itemize}

The complete numerical outputs include

\begin{itemize}[nosep]
\item \texttt{all\_structures\_metrics.csv};
\item \texttt{147mA\_all\_structures\_predictions.csv};
\item \texttt{training\_history.csv};
\item \texttt{data\_split\_blocks.csv};
\item the three trained model files.
\end{itemize}

\section{Summary}

The final simulation compares series, series--parallel, and parallel
identifiers using the same unseen 147~mA experiment. Cyclic non-overlapping blocks expose training, validation, and internal development testing to different chirp-frequency regions while guard gaps reduce window leakage.
The updated result set includes full-record and high-frequency comparisons,
separate measured-versus-predicted figures for each structure, direct error
plots, regression figures, and numerical accuracy metrics. These outputs make
the expected advantage of series--parallel one-step identification clear while
keeping parallel/free-running evaluation as the stricter autonomous test.

\begin{thebibliography}{9}

\bibitem{wang2017}
Y. Wang,
``A New Concept using LSTM Neural Networks for Dynamic System Identification,''
American Control Conference, 2017.

\bibitem{ogunmolu2016}
O. Ogunmolu, X. Gu, S. Jiang, and N. Gans,
``Nonlinear Systems Identification Using Deep Dynamic Neural Networks,''
arXiv:1610.01439, 2016.

\bibitem{bergmeir2018}
C. Bergmeir, R. J. Hyndman, and B. Koo,
``A Note on the Validity of Cross-Validation for Evaluating Autoregressive Time Series Prediction,''
Computational Statistics \& Data Analysis, 2018.

\bibitem{cerqueira2020}
V. Cerqueira, L. Torgo, and I. Mozeti\v{c},
``Evaluating Time Series Forecasting Models: An Empirical Study on Performance Estimation Methods,''
Machine Learning, 2020.

\end{thebibliography}

\end{document}
